%% Proofs for arbitrary finite-arity normal modal operators.

\section{Finitary modal Lindenbaum algebras and local completeness}

\begin{annexitem}{annex:p10.polyadic.1}{thm:p10-polyadic-algebra-completeness}{The finitary modal quotient}
\begin{proof}
\paragraph{Congruence for every operator.}
Classical propositional tautologies and modus ponens make
\(\equiv_L\) an equivalence relation compatible with the Boolean
connectives.  Suppose \(\varphi_i\equiv_L\psi_i\) for
\(1\leq i\leq n_f\).  Both implications between each pair are
theorems.  Coordinate monotonicity, with the other arguments fixed,
gives both implications between the formulas obtained by replacing
that one argument.  Propositional reasoning gives their equivalence.
Replace the arguments one at a time and use transitivity of
equivalence.  This proves
\[
 f(\varphi_1,\ldots,\varphi_n)
           \equiv_L f(\psi_1,\ldots,\psi_n).
\]
For a nullary symbol there are no arguments to compare: its quotient
interpretation is the single class \([f()]_L\).
Thus all quotient operations are well-defined.

The Boolean identities hold because their formula counterparts are
tautologies.  As in the unary case,
\[
 [\chi]_L=1\quad\Longleftrightarrow\quad\chi\in L,
 \qquad
 [\chi]_L=0\quad\Longleftrightarrow\quad\neg\chi\in L.
\]
Indeed, \(\chi\) is propositionally equivalent to
\(\chi\leftrightarrow\top\), and \(\neg\chi\) to
\(\chi\leftrightarrow\bot\).

\paragraph{Normality and additivity in the quotient.}
For any positive-arity coordinate, the normality axiom gives
\([f_i(\bot)]_L=0\).
To get equality in the additivity law, first use the displayed
additivity axiom.  The tautologies
\(\varphi\to\varphi\vee\psi\) and
\(\psi\to\varphi\vee\psi\), followed by coordinate monotonicity,
give
\[
 \vdash_L f_i(\varphi)\to f_i(\varphi\vee\psi),\qquad
 \vdash_L f_i(\psi)\to f_i(\varphi\vee\psi).
\]
Propositional reasoning combines them into the reverse implication.
Consequently
\[
 [f_i(\varphi\vee\psi)]_L
       =[f_i(\varphi)]_L\vee[f_i(\psi)]_L.
\]
Every finite tuple of quotient elements has formula representatives,
so these identities hold at every tuple of quotient arguments.
This requires only finitely many representatives for any one instance.
No such law is demanded of \([f()]_L\) when \(n_f=0\).
It follows that \(F_{\Lambda,L}\) is a \(\Lambda\)-BAO.

\paragraph{All valuations and theoremhood.}
For the canonical assignment \(\nu(p)=[p]_L\), induction on a
formula gives
\[
 \llbracket\chi\rrbracket_\nu^{F_{\Lambda,L}}=[\chi]_L.
\]
For an arbitrary valuation \(v:X\to F_{\Lambda,L}\) and a fixed
\(\lambda\in L\), choose representative formulas \(\rho_p\) only
for the finitely many variables occurring in \(\lambda\), with
\([\rho_p]_L=v(p)\).
Let \(\sigma(p)=\rho_p\) on that finite set and
\(\sigma(q)=q\) elsewhere.
Induction on the subformulas of \(\lambda\) proves
\[
 \llbracket\lambda\rrbracket_v^{F_{\Lambda,L}}
                  =[\sigma(\lambda)]_L.
\]
At an operator node the defining quotient operation applies to
the values of its finitely many immediate subformulas.
At a nullary node both sides are \([f()]_L\).
Uniform-substitution closure gives \(\sigma(\lambda)\in L\), so
the displayed value is \(1\).
Thus \(F_{\Lambda,L}\in\operatorname{Alg}_\Lambda(L)\).
There is no need to choose representatives for all classes or all
variables simultaneously.

If \(\varphi\notin L\), the canonical assignment gives a value
\([\varphi]_L\ne1\), providing the asserted algebraic counterexample.
Every member of \(L\) is valid in every member of
\(\operatorname{Alg}_\Lambda(L)\) by definition.
Together with the established quotient membership, these facts
prove the four equivalent conditions.

The rules also agree with algebraic validity directly.
Normality and additivity are the BAO equations; monotonicity in
each coordinate follows from additivity.
If \(\varphi\to\psi\) is valid, its values satisfy
\(\llbracket\varphi\rrbracket_v\leq\llbracket\psi\rrbracket_v\)
at each valuation, so applying an operator in one fixed coordinate
preserves the inequality.
Propositional tautologies, modus ponens, and substitution preserve
validity as well, the last by composition with formula evaluation.
Hence validating any additional axiom set validates its entire
normal \(\Lambda\)-logic closure.

If \(\bot\in L\), explosion gives every formula in \(L\), so all
classes coincide.  Conversely, a trivial quotient has
\([\bot]_L=1\), which gives \(\bot\in L\).
An algebra validating an inconsistent \(L\) must have \(0=1\),
and every trivial \(\Lambda\)-BAO validates all formulas.

\paragraph{Comparison with one unary diamond.}
For a unary \(\Diamond\), define \(\Box\chi=\neg\Diamond\neg\chi\).
The diamond normality and additivity equivalences, with Boolean
reasoning, give the theorems \(\Box\top\) and
\[
 \Box(\varphi\wedge\psi)
                 \leftrightarrow\Box\varphi\wedge\Box\psi.
\]
Diamond monotonicity, applied to the contrapositive of a theorem
\(\varphi\to\psi\), gives box monotonicity.
Apply that rule to the tautology
\(((\varphi\to\psi)\wedge\varphi)\to\psi\) and use the displayed
meet equivalence.  This proves the usual box distribution axiom.
If \(\chi\in L\), then \(\neg\chi\leftrightarrow\bot\) is a
theorem.  Operator congruence and diamond normality imply
\(\neg\Diamond\neg\chi\), so theorem necessitation is admissible.

Conversely, in the unary normal box presentation, necessitation and
distribution turn a theorem \(\varphi\to\psi\) into
\(\Box\varphi\to\Box\psi\).
They also prove \(\Box\top\) and the meet equivalence, as shown in
Annex~\ref{annex:p10.modal-algebra.1}.
Boolean duality therefore gives diamond monotonicity, normality,
and additivity.  The theorem closure conditions coincide.
Both local derivability relations use only premises, those theorems,
and modus ponens, so their local conventions coincide too.
\end{proof}
\end{annexitem}

\begin{annexitem}{annex:p10.polyadic.2}{thm:p10-polyadic-canonical-truth}{The descriptive canonical frame and truth lemma}
\begin{proof}
If \(L\) is consistent, the preceding result makes
\(F_{\Lambda,L}\) nontrivial.  Its filter \(\{1\}\) is proper,
and BPIT gives an ultrafilter.  Thus \(W_L\ne\varnothing\).
The construction of \(\mathcal D_L\) is exactly the descriptive
representation of \(F_{\Lambda,L}\) in
Theorem~\ref{thm:p10-bao-descriptive-representation}.
That theorem proves descriptiveness and identifies its admissible
algebra with \(F_{\Lambda,L}\) under the Stone isomorphism
\(a\mapsto\widehat a\).
In particular, the admissible sets are closed under every operator,
including the distinguished sets interpreting nullary symbols.

Let \(w:X\to\mathcal A_L\) be any admissible valuation.
Each \(w(p)\) has a unique preimage \(v(p)\in F_{\Lambda,L}\)
under the Stone map; these unique preimages define a valuation
without choosing syntactic representatives.
We prove by formula induction that
\[
 \{U:\mathcal D_L,w,U\models\chi\}
             =\widehat{\llbracket\chi\rrbracket_v^{F_{\Lambda,L}}}.
\]
The variable and Boolean cases use the Stone-map identities.
For a positive-arity operator \(f\), the relational satisfaction
clause is precisely its complex operation on the truth sets of
the arguments.  Apply the induction hypotheses to these sets and
then use Theorem~\ref{thm:p10-jonsson-tarski}:
\[
 f^{\mathcal D_L^+}(\widehat a_1,\ldots,\widehat a_n)
          =\widehat{f^{F_{\Lambda,L}}(a_1,\ldots,a_n)}.
\]
This uses the full simultaneous witness result already proved for
every finite arity.  For \(n_f=0\), the truth set is
\(\{U:R_f(U)\}=\widehat{[f()]_L}\) by definition, proving the
nullary step directly.

Since \(F_{\Lambda,L}\) validates every \(\lambda\in L\) at every
valuation, the displayed identity shows that \(\mathcal D_L\)
validates \(L\) under all admissible valuations.
Taking \(v=\nu\) yields
\[
 \mathcal D_L,v_0,U\models\chi
       \quad\Longleftrightarrow\quad[\chi]_L\in U.
\]
If \(\varphi\notin L\), then \(\neg[\varphi]_L\ne0\).
The principal filter above this nonzero element is proper;
extend it to an ultrafilter by BPIT.
The truth lemma refutes \(\varphi\) at that point under \(v_0\).
Thus validity on \(\mathcal D_L\) is exactly theoremhood.
\end{proof}
\end{annexitem}

\begin{annexitem}{annex:p10.polyadic.3}{thm:p10-polyadic-local-completeness}{Strong local descriptive-frame completeness}
\begin{proof}
\paragraph{Finite deduction.}
The local premise rules give
\[
 \Gamma\vdash_L\varphi
 \quad\Longleftrightarrow\quad
 \text{there is finite }\Gamma_0\subseteq\Gamma
 \text{ with }\vdash_L(\bigwedge\Gamma_0)\to\varphi.
\]
For the forward direction, collect the finitely many premises
used in the derivation and write \(\delta=\bigwedge\Gamma_0\).
Induct along the derivation to prove \(\delta\to\chi\in L\)
for each line \(\chi\).
For a premise line use its conjunction projection; for an
\(L\)-theorem use the tautology \(\chi\to(\delta\to\chi)\).
For a modus-ponens step use
\[
 (\delta\to(\alpha\to\beta))\to
             ((\delta\to\alpha)\to(\delta\to\beta)).
\]
There is no coordinate-monotonicity or substitution step on a
premise-dependent line.  Results of those theorem rules already
belong to \(L\).
Conversely, derive the finite conjunction from its premises and
use the displayed implication by modus ponens.
For \(\Gamma_0=\varnothing\), the conjunction is \(\top\).

\paragraph{Soundness at a point.}
Fix a frame validating \(L\), an admissible valuation, and a point
where every member of \(\Gamma\) is true.
All \(L\)-theorems are true there, and modus ponens preserves
truth there.  Induction over a local derivation proves soundness.
This argument applies to every general frame validating \(L\),
whether descriptive or not.

\paragraph{A descriptive counterexample to each failed derivation.}
Suppose \(\Gamma\not\vdash_L\varphi\).  Then \(L\) is consistent.
For every finite \(\Gamma_0\subseteq\Gamma\), finite deduction gives
\[
 \left(\bigwedge_{\gamma\in\Gamma_0}[\gamma]_L\right)
                                  \wedge\neg[\varphi]_L\ne0.
\]
If this meet were zero, Boolean order would give
\([\bigwedge\Gamma_0\to\varphi]_L=1\), and algebraic
completeness would put that implication in \(L\), a contradiction.
The same nonzero property implies that finite premise meets without
\(\neg[\varphi]_L\) are also nonzero.
Hence the set
\[
 S=\{[\gamma]_L:\gamma\in\Gamma\}\cup\{\neg[\varphi]_L\}
\]
generates a proper Boolean filter: its elements are those above a
finite meet from \(S\), and none of those meets is zero.
Apply BPIT to obtain an ultrafilter \(U\) containing this filter.
The canonical truth lemma makes every \(\gamma\in\Gamma\) true
at \(U\) under \(v_0\), and makes \(\varphi\) false there.
The preceding theorem places the nonempty descriptive frame
\(\mathcal D_L\) in \(\operatorname{DGen}_\Lambda(L)\).
It therefore supplies the required counterexample, even for an
arbitrary infinite premise set.
The same counterexample belongs to the larger class of all
nonempty general frames validating \(L\); together with soundness
this proves the additional assertion about that class.

If \(L\) is inconsistent, every formula belongs to \(L\) and is
derivable from every premise set.  No nonempty general frame
validates \(\bot\).  Both specified frame classes are therefore
empty, and their semantic consequences hold vacuously.
No empty pointed model is introduced.
\end{proof}
\end{annexitem}
